Clover algoritmus klastruje sekvencie DNA na základe ich počiatočných úsekov a usporadúva ich do stromovej štruktúry podobnej trie.
Trie funguje ako hierarchická usporiadaná stromová štruktúra, kde každý uzol reprezentuje spoločný prefix a vetvy vedú k sekvenciám, ktoré sa líšia v nasledujúcich pozíciách.

Nástroj Clover prijíma dáta vo formáte, v ktorom každý riadok obsahuje identifikátor a DNA sekvenciu, napríklad FASTA alebo TXT súbor s formátom \enquote{ID sekvencia}.
